\documentclass[11pt,oneside]{article}

\usepackage[letterpaper,margin=0.82in,headheight=15pt,footskip=30pt]{geometry}
\usepackage[T1]{fontenc}
\usepackage[utf8]{inputenc}
\usepackage{tgpagella}
\usepackage{tgheros}
\usepackage{inconsolata}

\usepackage{microtype}
\usepackage{amsmath,amssymb}
\usepackage{xcolor}
\usepackage{parskip}
\usepackage{setspace}
\usepackage{enumitem}
\usepackage{titlesec}
\usepackage{fancyhdr}
\usepackage{lastpage}
\usepackage{hyperref}
\usepackage{bookmark}
\usepackage[most]{tcolorbox}
\usepackage{ragged2e}
\usepackage{etoolbox}

\definecolor{ManifestInk}{HTML}{202124}
\definecolor{ManifestGray}{HTML}{5F6368}
\definecolor{ManifestLight}{HTML}{F3F4F5}
\definecolor{ManifestRule}{HTML}{8A8D91}

\hypersetup{
  colorlinks=true,
  linkcolor=ManifestInk,
  urlcolor=ManifestGray,
  pdftitle={The Ethical Quiet Quitting Manifesto},
  pdfauthor={Leah},
  pdfsubject={Repair first; withdraw honestly when repair fails},
  pdfkeywords={quiet quitting, ethics, work, labor, management, boundaries}
}

\setstretch{1.04}
\setlength{\parindent}{0pt}
\setlength{\parskip}{0.54em}
\setcounter{secnumdepth}{-1}
\setcounter{tocdepth}{2}
\emergencystretch=3em
\clubpenalty=10000
\widowpenalty=10000
\displaywidowpenalty=10000

\titleformat{\section}
  {\Large\sffamily\bfseries\color{ManifestInk}}
  {}{0pt}{}
  [\vspace{0.2em}\color{ManifestRule}\titlerule]
\titlespacing*{\section}{0pt}{2.2em}{0.8em}

\titleformat{\subsection}
  {\large\sffamily\bfseries\color{ManifestInk}}
  {}{0pt}{}
\titlespacing*{\subsection}{0pt}{1.5em}{0.45em}

\setlist[itemize]{leftmargin=1.45em,itemsep=0.22em,topsep=0.35em}
\setlist[enumerate]{leftmargin=1.65em,itemsep=0.32em,topsep=0.4em}
\providecommand{\tightlist}{%
  \setlength{\itemsep}{0.22em}\setlength{\parskip}{0pt}}

\let\oldquote\quote
\let\endoldquote\endquote
\renewenvironment{quote}
  {\begin{tcolorbox}[
      enhanced,
      breakable,
      colback=ManifestLight,
      colframe=ManifestRule,
      boxrule=0pt,
      leftrule=3pt,
      arc=0pt,
      outer arc=0pt,
      left=10pt,right=10pt,top=8pt,bottom=8pt,
      before skip=0.8em,after skip=0.8em
    ]\itshape}
  {\end{tcolorbox}}

\pagestyle{fancy}
\fancyhf{}
\fancyhead[L]{\sffamily\footnotesize\color{ManifestGray}The Ethical Quiet Quitting Manifesto}
\fancyhead[R]{\sffamily\footnotesize\color{ManifestGray}Leah}
\fancyfoot[C]{\sffamily\footnotesize\color{ManifestGray}\thepage\ / \pageref*{LastPage}}
\renewcommand{\headrulewidth}{0.3pt}
\renewcommand{\footrulewidth}{0pt}

\begin{document}

\begin{titlepage}
\thispagestyle{empty}
\vspace*{0.9in}

{\sffamily\bfseries\fontsize{31}{36}\selectfont\color{ManifestInk}
The Ethical Quiet Quitting\par
Manifesto\par}

\vspace{0.5em}
{\color{ManifestRule}\rule{\textwidth}{1.2pt}}
\vspace{1.0em}

{\sffamily\fontsize{17}{23}\selectfont\color{ManifestGray}
Repair First. Stop Hiding the Damage.\par
Withdraw Only When Repair Fails.\par}

\vfill

\begin{tcolorbox}[
  enhanced,
  colback=ManifestLight,
  colframe=ManifestRule,
  boxrule=0pt,
  leftrule=4pt,
  arc=0pt,
  left=14pt,right=14pt,top=13pt,bottom=13pt
]
{\Large\itshape We withdraw the subsidy, not the integrity.}
\end{tcolorbox}

\vfill

{\sffamily\large\bfseries Leah\par}
\vspace{0.2em}
{\sffamily\color{ManifestGray}August 14, 2026 \quad Version 1.0\par}
\vspace{0.55in}
\end{titlepage}

\pagenumbering{roman}
\thispagestyle{plain}
\tableofcontents
\clearpage
\pagenumbering{arabic}

\section{Preamble}\label{preamble}

Quiet quitting should not be the first move. It should be the last honest boundary after good-faith repair has failed.

Work is a cooperative system. A worker owes competent, safe, honest labor. An employer owes the conditions that make such labor possible: sufficient resources, intelligible priorities, usable feedback channels, and some plausible connection between effort and a self-directed future. When either side fails, the other may still owe an attempt at repair. But no worker owes unlimited self-erasure to keep a broken arrangement looking functional.

\begin{quote}
\textbf{The central rule:} Try to make the workplace better in every ethical way reasonably available. When those pathways have been tested and closed, stop using private sacrifice to counterfeit institutional health. Perform the real job safely, honestly, and sustainably. That is quiet quitting proper.
\end{quote}

This manifesto is not a defense of apathy. It is an ethic for people who care enough to try, care enough to tell the truth, and eventually care enough about themselves to stop feeding a system that converts every rescue into evidence that no rescue was needed.

\section{1. What Quiet Quitting Is, and What It Is Not}\label{what-quiet-quitting-is-and-what-it-is-not}

This manifesto uses \textbf{quiet quitting proper} in a precise sense:

\begin{quote}
Quiet quitting proper means competently performing the work one is paid and reasonably expected to perform, while ceasing the discretionary, uncompensated, unsustainable, or identity-consuming excess that has silently become part of the institution's operating model.
\end{quote}

It is not refusing to work. It is refusing to donate an undefined second job.

It is not:

\begin{itemize}
\tightlist
\item
  sabotaging output;
\item
  deliberately working below a reasonable professional standard;
\item
  falsifying records, hiding urgent risks, or manufacturing failures;
\item
  using food safety, customer welfare, product quality, or coworker hardship as leverage;
\item
  creating bottlenecks that would not otherwise exist;
\item
  retaliating against a decent direct supervisor for failures above their authority;
\item
  performing helplessness as a strategy;
\item
  confusing revenge with a boundary.
\end{itemize}

It is:

\begin{itemize}
\tightlist
\item
  working at a sustainable professional pace;
\item
  fulfilling the defined role without silently absorbing every missing role around it;
\item
  declining chronic unpaid availability and unbounded role expansion;
\item
  refusing to turn every managerial failure into a personal emergency;
\item
  making incompatible demands visible instead of privately reconciling them through exhaustion;
\item
  preserving energy, time, health, and agency for a life beyond the institution.
\end{itemize}

The ethical distinction is simple:

\[
\boxed{\text{Ceasing to subsidize a broken system} \neq \text{manufacturing damage to a system}.}
\]

\section{2. Why Workers Withdraw}\label{why-workers-withdraw}

People do not always quiet quit because they stopped caring. Often they stop caring only after care loses every actionable pathway.

The worker learns, repeatedly, that:

\[
\text{speaking} \nrightarrow \text{change},
\]

\[
\text{excellent work} \nrightarrow \text{meaningful advancement},
\]

\[
\text{identifying a problem} \nrightarrow \text{repair},
\]

\[
\text{sacrifice} \rightarrow \text{the new minimum expectation}.
\]

A promotion may exist in name but not in meaning. Requests for help may be acknowledged but not acted upon. The chain of command may transmit orders downward while filtering inconvenient reality upward. The employee may still possess the physical ability to act, but the routes connecting action to a self-directed future have narrowed or closed.

That is not merely low morale. It is a contraction of practical agency.

Yet perceived closure can be mistaken, temporary, or caused by an ordinary communication failure. That is why ethical quiet quitting begins with an obligation to test the closure rather than immediately assume it.

\section{3. Repair First: The Ethical Order of Operations}\label{repair-first-the-ethical-order-of-operations}

Quiet quitting becomes ethically serious only when it is the end of an ordered process rather than the beginning of an emotional reaction.

\subsection{3.1 Diagnose before accusing}\label{diagnose-before-accusing}

First determine what is actually happening.

Distinguish a structurally impossible workload from a temporarily difficult week. Separate understaffing from bad scheduling, missing equipment, poor sequencing, unclear priorities, training gaps, demand spikes, or simple misunderstanding. Notice where the real constraint lives and who actually has authority to change it.

Do not begin by assuming malice. Begin by reconstructing the system.

The question is not merely, ``Why am I suffering?'' It is:

\[
\text{What combination of demand, staffing, equipment, policy, and priority produces this outcome?}
\]

\subsection{3.2 Use the nearest functioning channel}\label{use-the-nearest-functioning-channel}

Start with the people who can still hear you.

Communicate with the direct supervisor. Explain the problem in concrete terms. Separate facts from interpretations. Describe the current resources, the demanded outcomes, the conflict between them, and the smallest realistic repair.

A good direct supervisor is not merely another rung in a hierarchy. They can become a local witness to reality, a translator across institutional cultures, and a source of provenance when higher levels would otherwise see only an unexplained bad number.

\subsection{3.3 Propose repair, not merely distress}\label{propose-repair-not-merely-distress}

Distress matters, but institutions often know how to ignore distress. They have more difficulty ignoring a clearly specified constraint paired with a feasible intervention.

Where possible, propose actions such as:

\begin{itemize}
\tightlist
\item
  adding labor at a specific bottleneck or time window;
\item
  changing a production sequence;
\item
  reducing or reordering lower-priority work;
\item
  repairing or replacing a capacity-limiting tool;
\item
  clarifying which standard takes precedence when standards conflict;
\item
  creating an escalation rule for predictable overload;
\item
  changing schedules to match observed demand rather than idealized demand.
\end{itemize}

The goal is not to do management's entire job for free. The goal is to ensure that a realistic pathway was offered before concluding that none exists.

\subsection{3.4 Force incompatible priorities into the open}\label{force-incompatible-priorities-into-the-open}

When the available resources cannot satisfy all demands, do not privately absorb the contradiction.

Say, in substance:

\begin{quote}
With the present staffing and demand, we can complete A and B safely, or A and C safely, but not A, B, and C at the demanded standard. Which takes priority?
\end{quote}

This is one of the cleanest ethical instruments available to a worker. It transfers the tradeoff to the level that owns the authority to make it.

When management chooses, follow the chosen priority honestly. The unchosen work is then not an unexplained personal failure. It is the visible consequence of an explicit allocation decision.

\subsection{3.5 Escalate proportionally}\label{escalate-proportionally}

Use the chain of command while it remains a communication channel. When it becomes only a shield against unwelcome information, use other ethical and available routes: documented escalation, established complaint systems, collective worker discussion, union representation where applicable, or external reporting for genuine safety or legal concerns.

Escalation should be truthful, proportional, and specific. It should aim at correction rather than humiliation.

\section{4. Communicate in a Channel the Institution Can Read}\label{communicate-in-a-channel-the-institution-can-read}

Workers often discover that explanations are treated as opinions while money, labor hours, waste, overtime, throughput, and customer complaints are treated as reality.

That does not mean money is the message. Money is only a channel.

A negative number by itself is radically ambiguous. It may be decoded as understaffing, incompetence, poor management, bad equipment, unusual demand, laziness, or random variation. A financial consequence without causal provenance can communicate the opposite of what the worker intended.

The useful signal is not:

\[
\Delta \$ < 0.
\]

It is:

\[
(\text{staffing},\ \text{demand},\ \text{equipment},\ \text{priority decision})
\longrightarrow
\text{observable outcome}.
\]

Every visible consequence should therefore remain coupled to its source-map.

Record what can honestly be recorded. Preserve staffing levels, workload, material shortages, equipment failures, explicit priority decisions, and the resulting backlog or cost. Do not falsify, exaggerate, or stage anything. The purpose is to make the real production function legible.

The principle is:

\begin{quote}
Do not create a failure in order to send a message. Stop erasing the evidence of a failure that already exists.
\end{quote}

\section{5. Ethical De-Buffering: Stop Counterfeiting Capacity}\label{ethical-de-buffering-stop-counterfeiting-capacity}

Many workplaces are not functioning at the capacity management believes they possess. They are functioning at that capacity plus an invisible worker subsidy.

Let

\[
Y = F(D,N,E,H),
\]

where:

\begin{itemize}
\tightlist
\item
  (Y) is observed output;
\item
  (D) is demand;
\item
  (N) is staffing;
\item
  (E) is equipment and process capacity;
\item
  (H) is hidden heroic compensation: unsustainable pace, role absorption, skipped recovery, exceptional flexibility, constant rescue, and unrecorded cognitive or emotional labor.
\end{itemize}

When workers repeatedly supply a large (H), management learns the wrong lesson. It observes acceptable output and concludes that (N) and (E) were sufficient. The worker's competence becomes evidence against the worker's own request for help.

\textbf{Ethical de-buffering} means returning (H) toward a sustainable baseline while holding safety, honesty, and professional integrity fixed.

It may mean:

\begin{itemize}
\tightlist
\item
  working quickly but not frantically;
\item
  taking agreed or required breaks;
\item
  declining to cover chronic vacancies as an indefinite personal obligation;
\item
  ending habitual off-clock problem-solving;
\item
  refusing to treat every predictable rush as an unforeseeable emergency;
\item
  allowing noncritical backlog to remain visible;
\item
  reporting unfinished lower-priority work instead of hiding it through unpaid effort;
\item
  requiring explicit authorization before taking on materially expanded duties;
\item
  preserving enough capacity to perform tomorrow's work rather than consuming tomorrow to save today.
\end{itemize}

Heroic effort may be freely chosen during a genuine emergency. It cannot ethically be converted into the permanent baseline and then used to prove that the emergency staffing level is adequate.

\begin{quote}
\textbf{The worker's body is not the slack variable in a staffing model.}
\end{quote}

Ethical de-buffering is not making the system fail. It is refusing to counterfeit success.

\section{6. Employer Size Does Not Settle the Ethics}\label{employer-size-does-not-settle-the-ethics}

A large conglomerate may be able to absorb a financial loss, yet redirect the immediate cost onto a local team, a good supervisor, or vulnerable coworkers. A small business may be genuinely fragile, yet use that fragility to claim an unlimited moral entitlement to unpaid labor.

Size matters, but it does not decide the case.

The morally relevant question is where the cascade lands.

Ask:

\begin{enumerate}
\def\labelenumi{\arabic{enumi}.}
\tightlist
\item
  Does the withdrawal expose an existing resource constraint, or create a new injury?
\item
  Will the institution absorb the consequence, or merely transfer it sideways onto coworkers?
\item
  Are customers, patients, guests, or the public being used as involuntary leverage?
\item
  Is a decent local manager being made the sole bearer of a failure created above their authority?
\item
  Is the withdrawal proportional to the hidden subsidy being removed?
\item
  Does the action improve causal legibility, or merely increase pain?
\end{enumerate}

The ethical aim is to move information upward without moving unnecessary suffering sideways.

A small business does not own a worker's life because it is fragile. A conglomerate does not become harmless because it is rich. The duty is to minimize unjust collateral harm while refusing indefinite self-sacrifice.

\section{7. Quiet Quitting Proper: The Last Boundary}\label{quiet-quitting-proper-the-last-boundary}

Quiet quitting proper becomes ethically justified when the following conditions substantially hold:

\begin{enumerate}
\def\labelenumi{\arabic{enumi}.}
\tightlist
\item
  \textbf{The problem is persistent and material.} It is not merely one disappointing shift or one denied preference.
\item
  \textbf{Good-faith repair has been attempted.} The worker has communicated, proposed, clarified, or escalated through reasonable available pathways.
\item
  \textbf{The pathways are closed, inert, or unreasonably costly.} Further attempts predictably produce no repair, retaliation, or more uncompensated labor without meaningful access to change.
\item
  \textbf{The excess effort is discretionary and structurally masking.} It exceeds the sustainable core role and allows under-resourcing to appear adequate.
\item
  \textbf{Core duties remain intact.} Safety, honesty, feasible quality, and professional conduct are preserved.
\item
  \textbf{The withdrawal is proportional.} The worker removes the hidden subsidy rather than creating an artificial deficit.
\item
  \textbf{Collateral harm is minimized.} Coworkers and customers are not treated as hostages in a conflict with management.
\item
  \textbf{The purpose is boundary and truth, not punishment.} Anger may be present, but revenge does not govern the design.
\end{enumerate}

At that point, the worker may ethically:

\begin{itemize}
\tightlist
\item
  perform the defined role competently and sustainably;
\item
  decline chronic role creep without compensation, authority, or agreement;
\item
  stop volunteering for predictable emergencies produced by normal planning;
\item
  maintain agreed availability rather than permanent accessibility;
\item
  refuse impossible combinations and request a priority decision;
\item
  stop taking personal ownership of outcomes determined by managerial resource choices;
\item
  conserve time and energy for health, relationships, education, art, organizing, job search, or another future;
\item
  remain courteous without manufacturing enthusiasm.
\end{itemize}

Quiet quitting proper is not the abandonment of standards. It is the narrowing of scope to the standards the institution actually resources.

\section{8. The Non-Negotiable Invariants}\label{the-non-negotiable-invariants}

No matter how closed the pathway becomes, the following remain fixed:

\subsection{Safety remains fixed}\label{safety-remains-fixed}

Do not compromise food safety, physical safety, sanitation, medical safety, public safety, or any other duty whose failure can seriously harm another person.

\subsection{Honesty remains fixed}\label{honesty-remains-fixed}

Do not falsify records, misstate capacity, invent incidents, conceal urgent defects, or manipulate numbers. A truthful signal is the entire point.

\subsection{Core professional integrity remains fixed}\label{core-professional-integrity-remains-fixed}

Do not intentionally produce bad work where competent work remains feasible within the real role and available resources. Scope may contract; integrity does not.

\subsection{Human beings are not leverage}\label{human-beings-are-not-leverage}

Do not use customers, patients, guests, coworkers, or a decent direct supervisor as expendable transmission media. Prefer effects that expose allocation decisions to effects that simply distribute misery.

\subsection{Proportionality remains fixed}\label{proportionality-remains-fixed}

Withdraw what was being improperly extracted. Do not add retaliation on top of the withdrawal.

\subsection{Accountability remains fixed}\label{accountability-remains-fixed}

A broken institution does not erase responsibility for one's own actions. Explanation is not exculpation. A boundary must still be owned.

\section{9. The Ethical Quiet Quitting Test}\label{the-ethical-quiet-quitting-test}

Before quiet quitting proper, ask:

\begin{itemize}
\tightlist
\item
  Have I identified the actual constraint rather than merely named my frustration?
\item
  Have I tried the nearest realistic repair pathway?
\item
  Have I communicated the problem in terms the receiver can act upon?
\item
  Have I proposed at least one feasible intervention or priority choice?
\item
  Have I preserved an honest causal record where possible?
\item
  Is the effort I plan to withdraw truly discretionary, uncompensated, unsustainable, or outside the reasonably defined role?
\item
  Will the withdrawal reveal an existing limit rather than manufacture a new failure?
\item
  Are safety, honesty, feasible quality, and coworker welfare protected?
\item
  Is the scale of withdrawal proportional to the hidden subsidy?
\item
  Am I trying to stop an extraction, or trying to make someone suffer?
\item
  Would I be willing to describe the action plainly to a fair-minded observer?
\item
  If management repaired the pathway in good faith, would I be willing to reconsider?
\end{itemize}

A ``no'' does not always forbid action. It identifies truth debt. Pay that debt where reasonably possible before treating withdrawal as ethically settled.

\section{10. The Worker's Pledge}\label{the-workers-pledge}

We will try to repair before we retreat.

We will tell the truth before we use silence.

We will distinguish a hard day from a structurally impossible demand.

We will not confuse endurance with loyalty, or exhaustion with excellence.

We will not make impossible workloads appear ordinary by consuming ourselves in private.

We will ask those with authority to choose among incompatible priorities.

We will preserve the connection between resource decisions and their consequences.

We will not sabotage.

We will not falsify.

We will not use safety, customers, or coworkers as leverage.

We will give honest labor for honest compensation.

We will help during genuine emergencies without allowing emergency effort to become the invisible permanent baseline.

We will withdraw the hidden subsidy, not our integrity.

And when every reasonable ethical pathway to improvement has closed, we reserve the right to quiet quit proper.

\section{11. A Message to Managers}\label{a-message-to-managers}

Your most capable employees may be hiding your worst process failures.

When a worker stops overfunctioning, the dysfunction may not be new. It may only have become visible. Do not ask first, ``How do we make this person care again?'' Ask:

\begin{quote}
What had this person been supplying that our official staffing, process, and compensation model did not contain?
\end{quote}

A hierarchy that carries commands downward but cannot carry unwelcome reality upward is not functioning as a communication system. It is functioning as a performance system in which each level edits the truth to look competent to the level above.

Repair the channel before demanding renewed investment.

Reward early warning. Permit explicit priority calls. Distinguish a worker's boundary from insubordination. Do not punish the person who makes a hidden constraint measurable. When an employee says that A, B, and C cannot all be completed safely with the available resources, management's job is to allocate, not to demand that the contradiction disappear inside the employee.

The fastest way to produce quiet quitting is to make voice causally irrelevant while continuing to demand emotional ownership.

\section{Conclusion: A Last Boundary, Not a First Weapon}\label{conclusion-a-last-boundary-not-a-first-weapon}

Quiet quitting should be a last ethical boundary, not a first weapon.

We repair first because work can matter, coworkers matter, customers matter, and institutions can sometimes learn. We diagnose, communicate, propose, document, prioritize, and escalate because withdrawal without inquiry can mistake a repairable problem for a closed world.

But repair is not a lifelong obligation to bleed into a closed circuit.

When words fail, records are ignored, priorities remain impossible, and the institution expects private sacrifice to bridge a structural gap it refuses to acknowledge, the worker may ethically stop paying that subsidy.

We do not manufacture failure. We stop concealing it.

We do not abandon duty. We define it.

We do not withdraw care from human beings. We withdraw unlimited institutional access to ourselves.

\textbf{Repair first. Make the constraints legible. Withdraw the hidden subsidy. Then, only when all else fails, quiet quit proper.}

\textbf{We withdraw the subsidy, not the integrity.}

\begin{center}\rule{0.5\linewidth}{0.5pt}\end{center}

\emph{Scope note:} This is an ethical framework, not legal advice. Employment contracts, collective bargaining agreements, workplace policies, professional duties, and applicable law may create additional obligations or protections.

\emph{Conceptual provenance:} Developed from Leah's TLICA work on agency, semantic interoperability, configurational action, and the separation of acts, outcomes, cascades, responsibility, and accountability, together with lived analysis of institutional work and chronic understaffing.

\end{document}
